% ============================================================================
% CHAPTER 4 SOLUTIONS GUIDE
% Exercise Solutions & Answer Keys
% ============================================================================
% Source: /markdown/ch4/C4AM_updated.md (Solutions Guide)
%         /markdown/ch4/Ch4AS_updated.md (Technical Exercise Solutions)
% Note: This file is for the Instructor's Edition, not the main book
% ============================================================================

\chapter{Solutions Guide: Chapter 4}
\label{ch:solutions-ch04}

\begin{chapterquote}{Solutions for exercises, activities, and discussion questions, numbered per level to match the teacher manual}
Framework-aligned answers enabling instructors to facilitate discussions on threshold decisions, error tradeoffs, and the value judgments embedded in HITL design. Discussion questions without a solution here are open-ended seminar prompts; target insights for every question appear in the teacher manual.
\end{chapterquote}

% ============================================================================
\section{Discussion Question Solutions}
% ============================================================================

\subsection{Introductory Level Solutions}

\subsubsection{Question 1: Smoke Detector Threshold}

\textbf{Prompt:} ``A smoke detector is set to a low threshold. What are the consequences? Would you set it differently for a hospital vs.\ a warehouse?''

\paragraph{Model Answer.}

\textbf{Consequences of a low threshold:}
\begin{itemize}
    \item \textit{Benefit:} Very high sensitivity --- almost never misses a real fire (low FN rate)
    \item \textit{Cost:} High false alarm rate --- cooking smoke, dust, humidity can trigger it
    \item False alarms train occupants to ignore the alarm (a form of alert fatigue), which is itself dangerous
\end{itemize}

\textbf{Hospital setting:}
\begin{itemize}
    \item High Type~II cost: patients who cannot self-evacuate quickly
    \item High Type~I cost: hospital evacuations are disruptive and can harm fragile patients
    \item Recommendation: slightly lower threshold, but invest heavily in HITL --- staff trained to quickly verify alarms before triggering full evacuation
    \item The HITL band here is the 30-second window where staff determine: real fire or false alarm before escalating
\end{itemize}

\textbf{Warehouse setting:}
\begin{itemize}
    \item Fewer occupants; typically healthier, mobile people
    \item Cost of FP: evacuation disruption, productivity loss
    \item Recommendation: standard or slightly higher threshold; false alarms are less costly
\end{itemize}

\textbf{Key learning:} The optimal threshold differs by context because the cost structure of the errors differs. This is a value judgment, not a technical one.

\paragraph{Grade Band Examples.}

\begin{description}
    \item[A-grade] ``The hospital faces asymmetric costs in both directions: missed fires are catastrophic for non-ambulatory patients, but false alarms are also harmful (staff cannot ignore an alarm in a hospital). The right design is a two-threshold system: a low threshold triggers a silent staff alert for verification; a confirmed alarm triggers the full evacuation protocol. This is a classic HITL band applied to physical safety.''

    \item[B-grade] Correctly identifies that hospitals have higher FN cost than warehouses; proposes different thresholds; may not identify the HITL band structure.

    \item[C-grade] Notes that the settings should differ without articulating the cost structure that drives the difference.
\end{description}

\subsubsection{Question 2: Job Application Threshold Change}

\textbf{Prompt:} ``A job application screening algorithm is set to `top 30\%.' Company becomes more selective. Should the threshold change? Who decides?''

\paragraph{Model Answer.}

\textbf{Should it change?} Yes. The threshold was set to route approximately 70\% of applications to rejection. If hiring strategy changes to a higher standard, the routing logic must change accordingly.

\textbf{The deeper issue:} Changing from ``top 30\%'' to ``top 15\%'' changes who bears the cost of each error type:
\begin{itemize}
    \item More conservative threshold: more good candidates rejected (FN cost falls on candidates)
    \item Less conservative threshold: more poor candidates sent to human review (FP cost falls on reviewers)
\end{itemize}

\textbf{Who decides?}
\begin{itemize}
    \item Not the algorithm alone --- it doesn't know what ``qualified'' means in the new strategy
    \item Not the data scientist alone --- they don't know the business priorities or legal constraints
    \item The people who bear the consequences: hiring managers (know job requirements), HR (know legal constraints), business leadership (know strategy)
    \item And critically: those who bear the cost of the errors (candidates), who are often not at the table
\end{itemize}

\paragraph{Common wrong answers.}
\begin{itemize}
    \item ``The algorithm handles it automatically'' --- No; the threshold is a design choice.
    \item ``The data scientist should optimize it'' --- Optimizing requires specifying a metric, which requires knowing whose interests to prioritize.
\end{itemize}

\subsubsection{Question 3: Probability Is Not Decision}

\textbf{Prompt:} ``What is the difference between `the model thinks this email is 70\% likely spam' and `this email is spam'? Why does it matter?''

\paragraph{Model Answer.}

\textbf{The difference:}
``70\% likely spam'' is a probability --- an expression of uncertainty about a specific email. ``This email is spam'' is a decision --- a binary commitment to action that discards the uncertainty information.

\textbf{Why it matters:}

\begin{enumerate}
    \item \textbf{Reversibility:} A probability can be acted upon at different thresholds depending on context. A label is applied uniformly and loses context.

    \item \textbf{Accountability:} If the email was legitimate and important, the label made the decision seem definitive. The probability flag would have preserved the option for reconsideration.

    \item \textbf{Threshold sensitivity:} The jump from 0.49 to 0.51 is statistically tiny but crosses the threshold and changes the decision entirely. This brittleness is a structural feature of any threshold system.

    \item \textbf{HITL opportunity:} A probability of 70\% is not 95\% --- it signals that this case might benefit from human review. Treating it identically to a 99\% spam email discards useful routing information.
\end{enumerate}

\subsubsection{Question 4: Medical Context --- ``Set the Threshold Conservatively''}

\textbf{Prompt:} ``A cancer screening AI has AUC = 0.95. The radiologist's supervisor says `set the threshold conservatively.' What does this mean in practice, and what tradeoff is it making?''

\paragraph{Model Answer.}

First, note that AUC says nothing about where to cut: AUC = 0.95 means the model ranks cases well, but the threshold is a separate, cost-driven decision.

``Conservatively'' in a screening context means \emph{do not miss cancers}: set the threshold low, so that almost all true cases are flagged (high sensitivity) at the price of more false positives --- healthy patients sent for follow-up imaging or biopsy.

\textbf{The tradeoff being made:} the supervisor is implicitly asserting that the cost of a missed cancer (delayed diagnosis, potentially fatal) far exceeds the cost of a false alarm (anxiety, an unnecessary procedure, expense). Formally, a low threshold is optimal when $C_{FN} \gg C_{FP}$.

\textbf{Worth surfacing in discussion:}
\begin{itemize}
    \item ``Conservative'' is itself a value judgment wearing a technical label --- someone decided whose costs count more.
    \item False positives are not free: over-testing causes real harm (biopsy complications, overtreatment of slow-growing cancers), so ``conservative'' can be contested by clinicians with a different cost view.
    \item The HITL framing: a low threshold paired with radiologist review of flagged cases is exactly a HITL band --- the model casts a wide net, humans make the costly call.
\end{itemize}

% ============================================================================
\subsection{Intermediate Level Solutions}
% ============================================================================

\subsubsection{Question 1: Hospital A vs.\ Hospital B Thresholds}

\textbf{Prompt:} ``AUC = 0.95. Hospital A sets $\tau = 0.3$, Hospital B sets $\tau = 0.6$. Describe expected outcomes.''

\paragraph{Model Answer.}

\textbf{Hospital A ($\tau = 0.3$ --- aggressive):}
\begin{itemize}
    \item True Positive Rate (sensitivity): very high ($\sim$95\%)
    \item False Positive Rate: high (30--40\% of healthy patients flagged)
    \item Consequences: many unnecessary biopsies, patient anxiety, healthcare costs, but very few missed cancers
\end{itemize}

\textbf{Hospital B ($\tau = 0.6$ --- conservative):}
\begin{itemize}
    \item True Positive Rate: moderate ($\sim$75\%)
    \item False Positive Rate: low ($\sim$8\%)
    \item Consequences: fewer unnecessary biopsies, but some cancers detected late
\end{itemize}

\textbf{Recommendation:} For cancer screening, the evidence strongly supports a lower threshold (Hospital A approach). Missing a cancer is typically irreversible and fatal; unnecessary biopsy is harmful but recoverable. Nuance: the optimal threshold also depends on cancer type --- for slow-growing conditions where overtreatment itself causes harm (certain prostate cancers), Hospital B's approach may be justified.

\textbf{HITL design implication:} Neither hospital should rely on a single automated threshold. The middle zone (0.3 to 0.6) should go to radiologist review. The hospitals' real choice is which side of the middle zone gets auto-approved and which gets auto-rejected.

\subsubsection{Question 2: Two-Threshold HITL System}

\textbf{Prompt:} ``Why is a HITL band (two thresholds) generally preferable to single-threshold for high-stakes applications?''

\paragraph{Model Answer.}

A single-threshold system makes a binary commitment at every point on the score scale. An email with score 0.51 goes to spam; an email with score 0.49 goes to inbox. The threshold is a single point of brittleness.

A two-threshold HITL band system:
\begin{itemize}
    \item Above $\tau_H$: high confidence positive $\to$ automated action
    \item Below $\tau_L$: high confidence negative $\to$ automated action
    \item Between $\tau_L$ and $\tau_H$: uncertain $\to$ human review
\end{itemize}

\textbf{Why this is better for high-stakes applications:}

\begin{enumerate}
    \item \textbf{Routes the highest-error-risk cases:} Cases near any single threshold are where errors are most likely. The band explicitly routes these to human review.

    \item \textbf{Separates the two error types:} $\tau_L$ and $\tau_H$ can be chosen independently, allowing separate optimization against the two error types.

    \item \textbf{Preserves precision at extremes:} The auto-zones achieve high precision on both sides; the uncertain middle zone gets expert attention.

    \item \textbf{Scales error cost appropriately:} The band can be placed asymmetrically around the expected optimal threshold --- a larger margin on the costly-error side.
\end{enumerate}

\textbf{Optimal band width is determined by:}
\begin{itemize}
    \item Cost of human review per case ($C_H$)
    \item Cost asymmetry ratio ($C_{FN}/C_{FP}$)
    \item Human reviewer accuracy advantage over model in the uncertain zone
    \item Human reviewer daily capacity
\end{itemize}

\subsubsection{Question 3: Allegheny Family Screening Tool}

\textbf{Prompt:} ``Critics argue the AFST threshold should be set by social workers with lived experience, not data scientists. Evaluate.''

\paragraph{Model Answer.}

\textbf{The critics are largely correct, and Chapter~4 explains why.}

\textbf{Why data scientists should not set the threshold alone:}
\begin{itemize}
    \item The threshold setting embeds a value judgment about the acceptable rate of traumatic false investigations vs.\ the acceptable rate of missed abuse
    \item That tradeoff is not answerable by the data --- it requires knowing what kind of society we want and who bears the cost of each error type
    \item Low-income families and families of color are disproportionately represented in the child welfare system; they bear asymmetric costs of false positives; data scientists are typically not from these communities
\end{itemize}

\textbf{Why social workers with lived experience should be involved:}
\begin{itemize}
    \item They understand the real cost of a false positive investigation (not just inconvenience --- potentially traumatic, with lasting consequences)
    \item They understand the heterogeneity in what a given risk score means for different family contexts
    \item They will implement the decisions; their calibration affects the system's real-world performance
\end{itemize}

\textbf{The appropriate process:}
\begin{itemize}
    \item Data scientists define the technical performance of the tool at various thresholds
    \item Domain experts (social workers, child welfare advocates, affected community representatives) deliberate on threshold placement given those tradeoffs
    \item The decision is made explicit, documented, and subject to review --- not embedded silently in a technical parameter
\end{itemize}

\begin{keyconcept}[Grading Note]
The strongest answers will recognize that this is not simply ``social workers are right, data scientists are wrong.'' It is about the process of threshold-setting: who participates, how value judgments are made explicit, and how accountability is maintained. A data scientist who defers the threshold to social workers without providing clear technical information about the tradeoffs has also failed their responsibility.
\end{keyconcept}

\subsubsection{Question 4: Dynamic Thresholds and Staffing}

\textbf{Prompt:} ``A fraud detection system has a HITL band calibrated to a team of 20 analysts. During holiday staffing, only 12 analysts are available. How should the band change? What are the consequences of each adjustment option?''

\paragraph{Model Answer.}

The band's width is a capacity decision: the review queue must not exceed what 12 analysts can clear. Three adjustment options:

\begin{enumerate}
    \item \textbf{Narrow the band symmetrically (raise $\tau_L$, lower $\tau_H$):} fewer cases routed to review; volume matches capacity. Consequence: more genuinely uncertain cases are auto-decided --- error rates rise on both sides, and the errors land in exactly the zone where the model is least reliable.
    \item \textbf{Narrow the band asymmetrically:} shrink the band more on the low-cost-error side and keep coverage on the high-cost side. Consequence: protects against the costly error type (e.g., missed fraud) while accepting more cheap errors (e.g., false declines) --- the cost-asymmetry logic of Chapter~4 applied to capacity.
    \item \textbf{Keep the band, let the queue grow (or triage within it):} cases wait, or reviewers spend less time per case. Consequence: delayed review of time-sensitive fraud, or shallower review --- the cost shows up as latency or as lower human accuracy instead of a routing change.
\end{enumerate}

\textbf{Key teaching point:} threshold changes are not only statistical decisions --- they are operational ones. Any band calibration that ignores staffing reality is a paper design; and a band silently narrowed for capacity is a value judgment made by circumstance rather than by decision-makers.

% ============================================================================
\subsection{Advanced Level Solutions}
% ============================================================================

\subsubsection{Question 1: Optimal Threshold Formula Derivation}

\textbf{Prompt:} ``Derive the optimal threshold formula $\tau^* = C_{FP} / (C_{FP} + C_{FN})$ from expected cost minimization, and explain why the class priors do not appear.''

\paragraph{Model Answer.}

Expected cost at threshold $\tau$:
\[
E[C(\tau)] = C_{FN} \cdot P(\hat{y}=0 \mid y=1) \cdot P(y=1) + C_{FP} \cdot P(\hat{y}=1 \mid y=0) \cdot P(y=0)
\]

For a calibrated classifier with score $s(x) = P(y=1 \mid x)$, the decision rule $\hat{y} = \mathbf{1}[s(x) \geq \tau]$ gives the boundary condition at score $\tau$:

At the boundary, choosing positive costs $C_{FP} \cdot P(y=0 \mid s=\tau)$ in expected terms; choosing negative costs $C_{FN} \cdot P(y=1 \mid s=\tau)$.

The optimal threshold is where these costs are equal:
\[
C_{FP} \cdot P(y=0 \mid s=\tau^*) = C_{FN} \cdot P(y=1 \mid s=\tau^*)
\]

Using the calibration assumption $P(y=1 \mid s=\tau) = \tau$ and Bayes theorem:
\[
C_{FP} \cdot (1 - \tau^*) = C_{FN} \cdot \tau^*
\]

Solving for $\tau^*$:
\[
\tau^* = \frac{C_{FP}}{C_{FP} + C_{FN}}
\]

Note that the class priors do not appear. This is not a balanced-classes assumption: for a calibrated score $s(x) = P(y=1 \mid x)$, the posterior already incorporates the priors, so $\tau^* = C_{FP}/(C_{FP} + C_{FN})$ holds at \emph{any} class balance (see Appendix~4A, Section on optimal threshold selection). Do not ``correct'' the formula by multiplying in the priors again --- that double-counts them. The prior-weighted expression
\[
\frac{s(x)}{1 - s(x)} \geq \frac{C_{FP} \cdot P(y=0)}{C_{FN} \cdot P(y=1)}
\]
is the \emph{likelihood-ratio} formulation of the same rule: it is the correct threshold when the score is a prior-free likelihood ratio rather than a calibrated posterior, not an adjustment to apply to posterior probabilities.

\textbf{As $C_{FN}/C_{FP} \to \infty$:} The denominator of $\tau^* = C_{FP}/(C_{FP} + C_{FN})$ is dominated by $C_{FN}$, and $\tau^* \to 0$. This means: flag every case as positive regardless of score. When missing a positive is infinitely more costly than a false alarm, the optimal strategy is to never miss. In practice, $C_{FN}/C_{FP}$ is large but finite, pushing $\tau^*$ to a very low but nonzero value --- consistent with aggressive cancer screening approaches.

% ============================================================================
\section{Activity Solutions}
% ============================================================================

\subsection{Activity 1: Cost Asymmetry Classification}

\begin{description}
    \item[FP more costly than FN] Antivirus software blocking legitimate software (FP: disabled critical app; FN: malware runs but may be caught later); parole denial for a safe person (FP: injustice, wasted incarceration cost; FN: reoffense, but many alternatives exist); voice authentication where re-authentication is easy (FP: denied access; FN: attacker waits for another chance).

    \item[FN more costly than FP] Cancer screening (FP: unnecessary biopsy; FN: missed cancer, potentially fatal); bridge safety inspection (FP: unnecessary inspection; FN: missed structural failure); wire transfer fraud detection (FP: delayed transfer; FN: money lost permanently).

    \item[Contested cases] These reflect values, not just technical preferences. The parole recommendation is a particularly rich example --- the answer depends on values about criminal justice.
\end{description}

\subsection{Activity 2: Confusion Matrix Exercise}

\paragraph{Given: 1,000 emails, 100 spam (10\% base rate). $\tau = 0.5$.}

\begin{table}[htbp]
\caption{Confusion matrix at $\tau = 0.5$}
\label{tab:conf-matrix-solution}
\centering
\begin{tabular}{@{}lll@{}}
\toprule
\textbf{Metric} & \textbf{Calculation} & \textbf{Value} \\
\midrule
Accuracy & $(80 + 850)/1000$ & 93.0\% \\
Precision & $80/(80+50)$ & 61.5\% \\
Recall (TPR) & $80/100$ & 80.0\% \\
FPR & $50/900$ & 5.6\% \\
\bottomrule
\end{tabular}
\end{table}

\paragraph{To achieve 95\% spam detection (Recall = 0.95):} Must lower $\tau$ until TPR = 0.95. This increases FPR substantially. The tradeoff: catching 15 more spam emails while blocking additional legitimate emails --- exactly how many cannot be computed from the $\tau = 0.5$ confusion matrix alone; the FPR at Recall = 0.95 depends on the full score distribution (the ROC curve), which is not given. Whether this tradeoff is acceptable is a value judgment --- and evaluating it quantitatively requires that curve.

% ============================================================================
\section{Common Wrong Answers}
% ============================================================================

\paragraph{``Just set the threshold to maximize accuracy.''}
Wrong. Accuracy is an average that masks the tradeoff. On an imbalanced problem (0.1\% fraud rate), a model that predicts ``not fraud'' for everything achieves 99.9\% accuracy while catching 0\% of fraud. Chapter~4's lesson is that accuracy is the wrong metric whenever error costs are asymmetric.

\paragraph{``The threshold doesn't matter much --- the model handles it.''}
Wrong. The threshold translates probability into action. Two systems with identical models but different thresholds have dramatically different false positive and false negative rates. The threshold is entirely a design choice.

\paragraph{``Higher threshold is always safer.''}
Wrong. ``Safer'' depends on which error you are trying to avoid. For cancer screening, a higher threshold means more missed cancers --- that is less safe for patients. For criminal sentencing, a higher recidivism-risk threshold means more restrictive release policies --- which may be less just.

\paragraph{``Setting the threshold is a purely technical decision.''}
This is Chapter~4's central corrective. The threshold embeds a value judgment about which error costs more. Calling it ``technical'' conceals the value judgment and removes it from deliberative accountability. A threshold that appears in code is no less a policy decision than a statute that appears in law.

% ============================================================================
\section{Technical Exercise Solutions (Appendix 4A)}
% ============================================================================

\subsection{ROC Curve Interpretation}

\paragraph{Key facts for instructors:}
\begin{itemize}
    \item AUC = 0.5: random classifier (diagonal ROC curve)
    \item AUC = 1.0: perfect classifier
    \item AUC = 0.95 is very good for most medical applications; real-world systems often achieve 0.75--0.90
    \item AUC is a threshold-independent summary: it averages over all possible thresholds
    \item AUC is insensitive to class imbalance but does not account for cost asymmetry --- a model can have AUC = 0.95 while being nearly useless at any practically achievable threshold
\end{itemize}

\subsection{HITL Band Optimization Example}

\textbf{Parameters:}
\begin{itemize}
    \item Review capacity: 20\% of cases
    \item $C_{FN} = 10$, $C_{FP} = 2$, $C_{\text{review}} = 1$
    \item Model AUC = 0.92
    \item Human reviewer accuracy in uncertain zone: 85\%
\end{itemize}

\textbf{Optimal band placement:}
\begin{enumerate}
    \item Compute optimal threshold without HITL: $\tau^* = 2/(2+10) \approx 0.17$
    \item At 20\% review capacity, the band should be placed around $\tau^*$ with width calibrated to capture exactly 20\% of cases
    \item Place $\tau_L$ and $\tau_H$ symmetrically if cost structure is symmetric; asymmetrically otherwise
    \item Given high $C_{FN}/C_{FP}$, the band should be placed lower (more conservative) --- slightly more cases go to the ``positive'' side of the auto-zone
\end{enumerate}

\textbf{Expected result:} Human review of the uncertain 20\% catches most of the false negatives while limiting false positives on the auto-positive end. Overall error cost is lower than either full automation or full human review.
